\section{Introduction}
\label{sec:introduction}

In \emph{Dynamics of Complex Systems}, Yaneer Bar-Yam describes \emph{emergent complexity}: a system composed of simple parts whose collective behaviour is complex~\cite{baryam1997}. The outcome of such a system cannot be determined by studying its parts on their own. Snowflakes are a familiar example: each grows outward from its centre by a small set of local rules, yet the order in which those rules act, and the way the result of one rule shapes the conditions for the next, produce dendrites that are never quite alike~\cite{snowflake2022}. A social group is a closer one. There are rules between each pair of people within it, and rules for the group as a whole, and how these interact under changing conditions is very hard to predict.

For artists, this has been a productive idea for more than half a century. Systems art, process art and generative art all take, in different ways, the position that an artist can author the conditions from which a work emerges rather than the work itself~\cite{burnham1968, galanter2003, boden2009}. In computational practice, the canonical demonstrations of emergence are cellular automata~\cite{wolfram2002, gardner1970}, flocking~\cite{reynolds1987}, and, closest to this work, Casey Reas's \emph{Process Compendium}, in which simple \emph{forms} (a circle, a line) are given simple \emph{behaviours} (move in a straight line, change direction when touching another element), and what is exhibited is not the elements but drawings made from their relationships~\cite{reas2010, reasgithub}.

These systems share a quality that this paper takes as its starting point: their rules are \emph{arbitrary}. They are elegant and productive, but they are chosen for the forms they generate, not because they are about anything in particular. A rule such as ``a cell with two equal neighbours turns on'' has no referent; it is a means. This is not a flaw in those works, which are about form and process. It does, however, leave open a question that motivates this one: what happens to emergence as an artistic material when the rules themselves are drawn from life?

This paper presents \emph{emergence-study}, an ongoing artwork and software system that takes that question literally. Its rules are drawn from the motions and physical constraints of everyday life, loosely inspired by the daily movement of people in New York. A \emph{world} contains \emph{places}. A \emph{population} of \emph{beings} is born, sustained and killed by the world over \emph{time}. Beings express life through \emph{movement}, which is governed by a \emph{schedule}: a daily routine of places they frequent. Both the schedule and the movement depend on \emph{age}; the closer a being is to its twenties, for example, the busier its day is likely to be. Every such rule carries noise, so that the unlikely (a busy sixty-year-old, a frail twenty-year-old) is possible but rare. The system is explicitly not representative of any real city. It is, as its documentation says, ``merely inspired by a tiny part of the multitude i / we live in.''

Crucially, the system is never shown. What a viewer sees is a \emph{software-interpretation}: a thought written in English, drawn from the artist's own experience, and then realised as an algorithm that selects data from the running system and represents it. One interpretation reads the system as threads of fate between beings born close to each other; another as ripples sent out when two beings meet; another as lines of separation drawn between neighbouring places. All of them run on the same world and the same rules.

\paragraph{Contributions.} This paper makes three contributions:
\begin{enumerate}[leftmargin=1.5em, itemsep=0.2em]
  \item \textbf{A base system of everyday rules} (Section~\ref{sec:system}): an agent-based world in which routines, ageing, energy, birth and death, rather than abstract local rules, are the source of emergence. It is described in full, with its equations, so that it can be rebuilt and reused.
  \item \textbf{The software-interpretation as a method} (Section~\ref{sec:method}): a three-part format (\emph{thought}, \emph{expression}, \emph{parameters}) through which an artist reads a running system. It situates software art within the lineage of instruction art~\cite{lewitt1967, ono1964}, while separating the reader of a system from its maker.
  \item \textbf{Six interpretations and a revision} (Sections~\ref{sec:interpretations} and~\ref{sec:revision}): six works made with the method, and a measured account of how revising the base system changed what any interpretation could show. The second is, I argue, a finding in its own right: the ground on which interpretations stand is not neutral.
\end{enumerate}

Section~\ref{sec:background} situates the work in the literatures of emergence, instruction and generative art, artificial life, and everyday life and mobility. Section~\ref{sec:discussion} discusses what everyday rules afford and what they do not, the authorship of interpretations, and the limits of the claims made here.

\anecdote{a short paragraph on where this began: the itp personal statement (2024), first year pulling in other directions, & finding the way back to this enquiry at the recurse center in summer 2026. reviewers at art venues respond well to a brief, honest origin; keep it to 3--4 sentences. (the quote from the statement could go here, or as an epigraph under the title.)}
